% \iffalse meta-comment
%
% Copyright (C) 2017 David Beauchemin, Vincent Goulet
%
% This file may be distributed and/or modified under the conditions
% of the LaTeX Project Public License, either version 1.3c of this
% license or (at your option) any later version. The latest version
% of this license is in:
%
%   http://www.latex-project.org/lppl.txt
%
% and version 1.3c or later is part of all distributions of LaTeX
% version 2006/05/20 or later.
%
% This work has the LPPL maintenance status `maintained'.
%
% The Current Maintainer of this work is Vincent Goulet
% <vincent.goulet@act.ulaval.ca>.
%
% This work consists of the files actuarialsymbol.dtx and the derived
% files.
%
% \fi
%
% \iffalse
%<*dtx>
          \ProvidesFile{actuarialsymbol.dtx}
%</dtx>
%<package>\NeedsTeXFormat{LaTeX2e}[2009/09/24]
%<package>\ProvidesPackage{actuarialsymbol}[2017/04/10 v0.1 Actuarial symbols]
%<*batchfile>
\begingroup
\input docstrip.tex
\keepsilent
\askforoverwritefalse

\preamble

This is a generated file.

Copyright (C) 2017 David Beauchemin, Vincent Goulet

This file may be distributed and/or modified under the conditions
of the LaTeX Project Public License, either version 1.3c of this
license or (at your option) any later version. The latest version
of this license is in:

   http://www.latex-project.org/lppl.txt

and version 1.3c or later is part of all distributions of LaTeX
version 2006/05/20 or later.

This work has the LPPL maintenance status `maintained'.

The Current Maintainer of this work is Vincent Goulet
<vincent.goulet@act.ulaval.ca>.

This work consists of actuarialsymbol.dtx and the derived files.

\endpreamble
\generate{\file{actuarialsymbol.sty}{\from{actuarialsymbol.dtx}{package}}}
\endgroup
%</batchfile>
%<*driver>
\documentclass[11pt,x11names]{ltxdoc}
  \usepackage[utf8]{inputenc}
  \usepackage[T1]{fontenc}
  \usepackage{natbib}
  \usepackage{microtype}
  \usepackage[scaled=0.90]{helvet}
  \usepackage[scaled=1.02]{inconsolata}
  \usepackage[sc]{mathpazo}
  \usepackage{rsfso}
  \usepackage{actuarialsymbol}
  \usepackage{graphicx,framed}
  \usepackage{longtable,booktabs}
  \DisableCrossrefs
  \CodelineNumbered
  \RecordChanges
  \GlossaryPrologue{\section*{Version history}%
    \addcontentsline{toc}{section}{Version history}}

  %% colors
  \usepackage{xcolor}
  \definecolor{link}{rgb}{0,0.4,0.6}   % ~dvips RoyalBlue
  \definecolor{url}{rgb}{0.6,0,0}      % dark red
  \definecolor{citation}{rgb}{0,0.5,0} % dark green
  \colorlet{shadecolor}{LightYellow1}

  %% left aligned description list
  \usepackage{enumitem}
  \setlist[description]{leftmargin=*,align=left}

  %% environment for remarks
  \usepackage{amsthm}
  \theoremstyle{definition}
  \newtheorem*{rem}{Remark}

  %% single entry for symbol and shortcut lists
  %% http://tex.stackexchange.com/a/128441
  \usepackage{xparse}
  \ExplSyntaxOn
  \NewDocumentCommand{\showcase}{v} { \tl_rescan:nn { } { $#1$ } & \texttt{#1} }
  \NewDocumentCommand{\rshowcase}{v} { \texttt{#1} & \tl_rescan:nn { } { $#1$ } }
  \ExplSyntaxOff

  %% environments for demo code; 'texample' mostly lifted from
  %% documentation of amsmath
  \newenvironment{demo}{\begin{trivlist}\item}{\end{trivlist}}
  \newenvironment{texample}[1][0.5]{%
    \noindent\begin{minipage}{#1\linewidth}\begin{snugshade*}%
      \def\producing{\end{snugshade*}\end{minipage}\hfill\begin{minipage}{\dimexpr0.97\linewidth-#1\linewidth}%
        \hbox\bgroup\kern-.2pt%
        \vbox\bgroup\parindent0pt\relax
        % The 3pt is to cancel the -\lineskip from \displ@y
        \abovedisplayskip3pt \abovedisplayshortskip\abovedisplayskip
        \belowdisplayskip0pt \belowdisplayshortskip\belowdisplayskip
        \noindent}
    }{%
      \par
      % Ensure that a lonely \[\] structure doesn't take up width less than
      % \hsize.
      \hrule height0pt width\hsize
      \egroup\kern-.2pt\egroup
    \end{minipage}%
    \par
  }

  %% utility macro
  \newcommand{\pkg}[1]{\textsf{#1}}

  %% increase row height in tables
  \renewcommand{\arraystretch}{1.2}

  %% hyperlinks
  \usepackage{hyperref}
  \hypersetup{colorlinks, linktocpage,
    urlcolor=url, linkcolor=link, citecolor=citation,
    bookmarksopen, bookmarksnumbered, bookmarksdepth=subsection,
    pdftitle={actuarialsymbol User Guide},
    pdfauthor={David Beauchemin, Vincent Goulet}}

  %% autoref labels
  \let\subsectionautorefname\sectionautorefname
  \let\subsubsectionautorefname\sectionautorefname

\begin{document}
  \DocInput{actuarialsymbol.dtx}
\end{document}
%</driver>
% \fi
% \CheckSum{531}
%
% \GetFileInfo{actuarialsymbol.sty}
%
% \title{Actuarial symbols of life contingencies \\
%   and financial mathematics%
%   \thanks{This document corresponds to
%   \pkg{actuarialsymbol}~\fileversion, dated~\filedate.}}
% \author{David Beauchemin \\
%   \url{david.beauchemin.5@ulaval.ca}
%   \and
%   Vincent Goulet\thanks{%
%     École d'actuariat, Université Laval} \\
%   \url{vincent.goulet@act.ulaval.ca}}
% \maketitle
%
% \begin{abstract}
%   Package \pkg{actuarialsymbol} provides facilities to compose
%   actuarial symbols of life contingencies and financial mathematics
%   characterized by subscripts and superscripts on both sides of a
%   principal symbol. The package also features commands to easily and
%   consistently position precedence numbers above or below statuses
%   in symbols for multiple lives contracts.
%
%   Since the actuarial notation can get quite involved. the package
%   defines a number of shortcut macros to ease entry of the most
%   common elements. \autoref{sec:comprehensive} lists the commands to
%   typeset a large selection of symbols of life contingencies.
% \end{abstract}
%
% \section{Introduction}
% \label{sec:introduction}
%
% Actuaries denote various quantities of life contingencies like
% present values of life insurances and life annuities, annual
% premiums, or reserves using a whole array of symbols. The highly
% descriptive, yet compact, notation was standardized as far back as
% in 1898 \citep{Wolthuis:notation:2004}. \autoref{fig:mosaic} shows a
% creative use of the notation by the graduating class of 1972 in
% Actuarial Science at Université Laval.
%
% \begin{figure}[t]
%   \centering
%   \includegraphics[width=\linewidth]{mosaic}
%   \caption{``Actuariat'' (French for Actuarial Science) written
%     using actuarial symbols on the 1972 graduating class mosaic at
%     Université Laval}
%   \label{fig:mosaic}
% \end{figure}
%
% \citet[Appendix~4]{Bowers:2e:1997} offer an excellent overview of
% the composition rules for symbols of actuarial functions. In a
% nutshell, a principal symbol, say $S$, is combined with auxiliary
% symbols positioned in subscript or in superscript, to the left or to
% the right. Schematically, we thus have:
% \begin{equation}
%   \label{schema1}
%   \actsymb[\fbox{\scriptsize I}\!\;][\fbox{\scriptsize II}\!\;]%
%      {\framebox[5ex]{$S$}}{\,\fbox{\scriptsize III}}%
%      [\,\fbox{\scriptsize IV}]
% \end{equation}
%
% The principal symbol is in general a single letter. The letter may
% be ``accented'' with a bar ($\bar{A}$), double dots ($\ddot{a}$) or
% a circle ($\mathring{e}$). When the principal symbol consists of two
% letters, they are grouped between parentheses, as in $\IA$ or
% $\DA*$. Most commonly, there are alphanumeric statuses in the
% lower-right position \fbox{\small III}. Numerals can be placed above
% or below the individual statuses to show the order of failure; we
% will refer to these numerals as \emph{precedence numbers}. Otherwise,
% auxiliary symbols appear lower-left \fbox{\small I}, upper-left
% \fbox{\small II} and upper-right \fbox{\small IV}, in that order of
% frequency.
%
% Principle symbols for benefit premiums, reserves and amount of
% reduced paid-up insurance, $P$, $V$ and $W$, are combined with
% benefit symbols unless the benefit is a level unit insurance payable
% at the end of the year of death. In such cases, we have the
% following symbol structure (replace $P$ by $V$ or $W$ as needed):
% \begin{equation}
%   \label{schema2}
%   \actsymb[\fbox{\scriptsize I}\!\;][\fbox{\scriptsize II}\!\;]%
%      {\framebox[5ex]{$P$}}{}[\,\fbox{\scriptsize IV}]%
%   \bigl(\actsymb{\framebox[5ex]{$S$}}{\,\fbox{\scriptsize III}}\bigr)%
% \end{equation}
%
% Perhaps the most commonly used auxiliary symbol not readily
% available in {\LaTeX} is the ``angle'' denoting a duration $n$, as
% in $\angln$. Package \pkg{actuarialangle} \citep{actuarialangle}
% provides commands to create this symbol. This package is imported at
% load time by \pkg{actuarialsymbol}.
%
% Package \pkg{actuarialsymbol} provides a generic command to position
% all subscripts and superscripts easily and consistently around a
% principal symbol, a command to create two-letter symbols and two
% commands to position precedence numbers above and below statuses.
% Since entering actuarial symbols can get quite involved, the package
% also defines a number of shortcuts to create the most common
% actuarial functions of financial mathematics and life contingencies.
%
% \section{For the impatient}
% \label{sec:actuarialsymbol:quickstart}
%
% We provide in \autoref{sec:comprehensive} a comprehensive list of
% symbols of life contingencies along with the {\LaTeX} code to
% compose them with \pkg{actuarialsymbol} loaded. This may be all the
% hurried reader needs. In our highly biased view, it remains worth
% going through the fine documentation below.
%
% \section{Package features}
% \label{sec:features}
%
% This section describes the generic commands provided by the package
% to compose actuarial symbols. On a daily basis, one should use the
% shortcuts of \autoref{sec:shortcuts} to ease entry of symbols.
%
% \subsection{Actuarial symbol}
% \label{sec:features:actsymb}
%
% \begin{DescribeMacro}{\actsymb}
%   The generic command \cmd{\actsymb} typesets a principal symbol with
%   surrounding subscripts and superscripts. Its syntax is somewhat
%   unusual for {\LaTeX}, but it serves well the natural order of the
%   building blocks of a symbol and their relative prevalence:
%   \begin{quote}
%     \cmd{\actsymb}\oarg{ll}\oarg{ul}\marg{symbol}\marg{lr}\oarg{ur}
%   \end{quote}
%   Above, \meta{ll} identifies the auxiliary symbol in the lower left
%   subscript position \fbox{\small I} (following the notation in the
%   schematic representation \eqref{schema1}); \meta{ul} is the upper
%   left superscript \fbox{\small II}; \meta{symbol} is the principal
%   symbol $S$; \meta{lr} is the lower right subscript \fbox{\small
%   III}; \meta{ur} is the upper right superscript $\fbox{\small IV}$.
%   The principal symbol and the right subscript are required, the
%   other arguments are optional.
%   \begin{demo}
%     \begin{texample}
%       |\actsymb{A}{x}|
%       \producing
%       $\actsymb{A}{x}$
%     \end{texample}
%     \begin{texample}
%       |\actsymb[n]{A}{x}|
%       \producing
%       $\actsymb[n]{A}{x}$
%     \end{texample}
%     \begin{texample}
%       |\actsymb[n][2]{A}{x}|
%       \producing
%       $\actsymb[n][2]{A}{x}$
%     \end{texample}
%     \begin{texample}
%       |\actsymb[n][2]{A}{x}[(m)]|
%       \producing
%       $\actsymb[n][2]{A}{x}[(m)]$
%     \end{texample}
%   \end{demo}
%
%   The command actually admits one more optional argument to compose
%   symbols for premiums, reserves and paid-up insurance. The extended
%   command
%   \begin{quote}
%     \cmd{\actsymb}\oarg{ll}\oarg{ul}\oarg{principle}\marg{symbol}\marg{lr}\oarg{ur}
%   \end{quote}
%   puts symbol \meta{principle} outside the parentheses in the
%   schematic representation \eqref{schema2}.
%   \begin{demo}
%     \begin{texample}[0.7]
%       |\actsymb[][][P]{\bar{A}}{x:\angln}|
%       \producing
%       $\actsymb[][][P]{\bar{A}}{x:\angln}$
%     \end{texample}
%     \begin{texample}[0.7]
%       |\actsymb[k][][V]{\bar{A}}{x}[\{1\}]|
%       \producing
%       $\actsymb[k][][V]{\bar{A}}{x}[\{1\}]$
%     \end{texample}
%     \begin{texample}[0.7]
%       |\actsymb[k][][\bar{W}]{\bar{A}}{x}|
%       \producing
%       $\actsymb[k][][\bar{W}]{\bar{A}}{x}$
%     \end{texample}
%   \end{demo}
% \end{DescribeMacro}
%
% \begin{rem}
%   {\TeX} adjusts the position of a subscript downward when a
%   superscript is present:
%   \begin{displaymath}
%     A_x \quad A_x^2.
%   \end{displaymath}
%   The package maintains this behaviour. Therefore, entering the
%   symbols above using the standard operators |^| and |_| or with
%   \cmd{\actsymb} yields the same result.
%   \begin{demo}
%     \begin{texample}[0.7]
%       |A_x \quad A_x^2|
%       \producing
%       $A_x \quad A_x^2$
%     \end{texample}
%     \begin{texample}[0.7]
%       |\actsymb{A}{x} \quad \actsymb{A}{x}[2]|
%       \producing
%       $\actsymb{A}{x} \quad \actsymb{A}{x}[2]$
%     \end{texample}
%   \end{demo}
%   However, the package also ensures that the left and right
%   subscripts, when both present, are at the same level, something
%   popular ad~hoc constructions do not provide.
%   \begin{demo}
%     \begin{texample}[0.7]
%       |{}_tA_x \quad {}_tA_x^2|
%       \producing
%       ${}_tA_x \quad {}_tA_x^2$
%     \end{texample}
%     \begin{texample}[0.7]
%       |\actsymb[t]{A}{x} \quad \actsymb[t]{A}{x}[2]|
%       \producing
%       $\actsymb[t]{A}{x} \quad \actsymb[t]{A}{x}[2]$
%     \end{texample}
%   \end{demo}
%   Authors who would prefer a uniform subscript position \emph{throughout
%   their document} should load package \pkg{subdepth} \citep{subdepth}.
% \end{rem}
%
% \subsection{Two-letter principal symbols}
% \label{sec:features:twoletsymb}
%
% \begin{DescribeMacro}{\twoletsymb}
%   Entering two-letter principal symbols like $\DA$ as |$(DA)$| results in
%   letters that are too distant from one another: $(DA)$. To unify
%   presentation, the package provides the command
%   \begin{quote}
%     \cmd{\twoletsymb}\oarg{length}\marg{symbol\_1}\marg{symbol\_2}
%   \end{quote}
%   to group \meta{symbol\_1} and \meta{symbol\_2} between parentheses
%   with kerning\footnote{Spacing adjustment between the characters}
%   reduced by length \cmd{\twoletkern} (see below). One can also
%   reduce spacing by \meta{length} for a specific symbol.
%   \begin{demo}
%     \begin{texample}
%       |\twoletsymb{\bar{D}}{\bar{A}}|
%       \producing
%       $\twoletsymb{\bar{D}}{\bar{A}}$
%     \end{texample}
%     \begin{texample}
%       |\twoletsymb{I}{\ddot{a}}|
%       \producing
%       $\twoletsymb{I}{\ddot{a}}$
%     \end{texample}
%     \begin{texample}
%       |\twoletsymb[0.8pt]{I}{\ddot{a}}|
%       \producing
%       $\twoletsymb[0.8pt]{I}{\ddot{a}}$
%     \end{texample}
%   \end{demo}
%
%   We expect authors to use \cmd{\twoletsymb} to define commands, not
%   directly in equations. The package already defines a number of
%   shortcuts for the main two-letter actuarial symbols; see
%   \autoref{sec:shortcuts}.
% \end{DescribeMacro}
%
% \begin{DescribeMacro}{\twoletkern}
%   The standard kerning between mathematical symbols defined with
%   \cmd{\twoletsymb} is \emph{reduced} by the length
%   \cmd{\twoletkern}, by default \the\twoletkern. This value can be
%   changed as usual using \cmd{\setlength}.
% \end{DescribeMacro}
%
% \subsection{Precedence numbers}
% \label{sec:features:precedence}
%
% \begin{DescribeMacro}{\nthtop}
% \begin{DescribeMacro}{\nthtop*}
%   Precedence numbers appear above or below individual statuses in
%   the right subscript \fbox{\small III} of a symbol. Commands
%   \begin{quote}
%     \cmd{\nthtop}\oarg{length}\marg{number}\marg{status} \\
%     \cmd{\nthtop*}\oarg{length}\marg{number}\marg{status}
%   \end{quote}
%   put a precedence \meta{number} above a \meta{status}, smashed so
%   that the apparent height of the status is its normal height. This is
%   normally used in the right subscript \fbox{\small III} of a symbol. With
%   \cmd{\nthtop}, the spacing between the precedence number and the
%   status is a constant \cmd{\nthtopsep} (see below). This can result
%   in precedence numbers placed at different heights if one status
%   contains an horizontal rule.
%   \begin{demo}
%     \begin{texample}[0.85]
%       |\actsymb{A}{\nthtop{1}{x}:\angln}|
%       \producing
%       \rule{0pt}{12pt}$\actsymb{A}{\nthtop{1}{x}:\angln}$
%     \end{texample}
%     \begin{texample}[0.85]
%       |\actsymb{A}{x:\nthtop{1}{\angln}}|
%       \producing
%       \rule{0pt}{12pt}$\actsymb{A}{x:\nthtop{1}{\angln}}$
%     \end{texample}
%     \begin{texample}[0.85]
%       |\actsymb{A}{\nthtop{1}{x}y:\nthtop{2}{\angln}}|
%       \producing
%       \rule{0pt}{12pt}$\actsymb{A}{\nthtop{1}{x}y:\nthtop{2}{\angln}}$
%     \end{texample}
%   \end{demo}
%   Conversely, \cmd{\nthtop*} always leaves enough space
%   \cmd{\nthtopskip} for intervening horizontal rules, resulting
%   in vertically aligned precedence numbers.
%   \begin{demo}
%     \begin{texample}[0.85]
%       |\actsymb{A}{\nthtop*{1}{x}:\angln}|
%       \producing
%       \rule{0pt}{12pt}$\actsymb{A}{\nthtop*{1}{x}:\angln}$
%     \end{texample}
%     \begin{texample}[0.85]
%       |\actsymb{A}{x:\nthtop*{1}{\angln}}|
%       \producing
%       \rule{0pt}{12pt}$\actsymb{A}{x:\nthtop*{1}{\angln}}$
%     \end{texample}
%     \begin{texample}[0.85]
%       |\actsymb{A}{\nthtop*{1}{x}y:\nthtop*{2}{\angln}}|
%       \producing
%       \rule{0pt}{12pt}$\actsymb{A}{\nthtop*{1}{x}y:\nthtop*{2}{\angln}}$
%     \end{texample}
%   \end{demo}
%   The optional argument \meta{length} changes the default spacing
%   for one symbol. The package also defines shortcuts for first,
%   second and third top precedence; see \autoref{sec:shortcuts}.
% \end{DescribeMacro}
% \end{DescribeMacro}
%
% \begin{DescribeMacro}{\nthbottom}
% \begin{DescribeMacro}{\nthbottom*}
%   In the same vein as the above two commands,
%   \begin{quote}
%     \cmd{\nthbottom}\oarg{length}\marg{number}\marg{status} \\
%     \cmd{\nthbottom*}\oarg{length}\marg{number}\marg{status}
%   \end{quote}
%   put a precedence \meta{number} below a \meta{status}. With
%   \cmd{\nthbottom} the spacing between the status and the number is
%   a constant \cmd{\nthbottomsep}, whereas with \cmd{\nthbottom*}
%   enough spacing \cmd{\nthbottomskip} is left to bottom align the
%   precedence numbers. The |*| variant is most useful in cases
%   involving more than one bottom precedence numbers and statuses
%   with descenders.
%   \begin{demo}
%     \begin{texample}[0.85]
%       |\actsymb{A}{w\nthtop{3}{x}%| \\
%       |            \nthbottom{1}{y}\nthbottom{2}{z}}|
%       \producing
%       \rule{0pt}{12pt}$\actsymb{A}{w\nthtop{3}{x}\nthbottom{1}{y}\nthbottom{2}{z}}$
%     \end{texample}
%     \begin{texample}[0.85]
%       |\actsymb{A}{w\nthtop{3}{x}%| \\
%       |            \nthbottom*{1}{y}\nthbottom*{2}{z}}|
%       \producing
%       \rule{0pt}{12pt}$\actsymb{A}{w\nthtop{3}{x}\nthbottom*{1}{y}\nthbottom*{2}{z}}$
%     \end{texample}
%   \end{demo}
%   The optional argument \meta{length} changes the default spacing
%   for one symbol. The package also defines shortcuts for first,
%   second and third bottom precedence; see \autoref{sec:shortcuts}.
% \end{DescribeMacro}
% \end{DescribeMacro}
%
% \begin{rem}
%   The fact that top precedence numbers have zero height means they
%   will clash with a right superscript \fbox{\small IV}:
%   \begin{demo}
%     \begin{texample}[0.7]
%       |\actsymb{A}{\nthtop{1}{x}:\angln}[(m)]|
%       \producing
%       $\actsymb{A}{\nthtop{1}{x}:\angln}[(m)]$
%     \end{texample}
%   \end{demo}
%   In such rare circumstances, one needs to insert a \emph{strut} (an
%   invisible vertical rule) in the subscript to push it downward as
%   needed:
%   \begin{demo}
%     \begin{texample}[0.7]
%       |\actsymb{A}{\rule{0pt}{2.3ex}%| \\
%       |            \nthtop{1}{x}:\angln}[(m)]|
%       \producing
%       $\actsymb{A}{\rule{0pt}{2.3ex}\nthtop{1}{x}:\angln}[(m)]$
%     \end{texample}
%   \end{demo}
%   This remark also applies to bottom precedence numbers in inline
%   formulas or multiline equations.
% \end{rem}
%
% \begin{DescribeMacro}{\nthtopsep}
% \begin{DescribeMacro}{\nthtopskip}
% \begin{DescribeMacro}{\nthbottomsep}
% \begin{DescribeMacro}{\nthbottomskip}
%   The constant spacing between a top precedence number and the status
%   underneath when using \cmd{\nthtop} is \cmd{\nthtopsep}, by
%   default \the\nthtopsep. The constant height of top precedence
%   numbers when using \cmd{\nthtop*} is achieved by setting the
%   baseline skip to \cmd{\nthtopskip}, by default
%   \the\nthtopskip.
%
%   Similarly, the constant spacing between a bottom precedence number
%   and the status above when using \cmd{\nthbottom} is
%   \cmd{\nthbottomsep}, by default \the\nthbottomsep, and the
%   constant height of bottom precedence numbers when using
%   \cmd{\nthbottom*} is achieved by setting the baseline skip to
%   \cmd{\nthbottomskip}, by default \the\nthbottomskip.
%
%   These values can be changed as usual using \cmd{\setlength}.
% \end{DescribeMacro}
% \end{DescribeMacro}
% \end{DescribeMacro}
% \end{DescribeMacro}
%
% \section{Shortcuts}
% \label{sec:shortcuts}
%
% Composing actuarial symbols from scratch using \cmd{\actsymb} can
% easily get quite involved. For this reason, the package defines a
% large number of shortcut macros to ease entry of the most common
% symbols. We encourage authors to define their own shortcuts for
% cases we did not consider. The package provides facilities to define
% powerful shortcuts.
%
% \subsection{Basic symbols of life tables, insurance and annuities}
% \label{sec:shortcuts:principal}
%
% \autoref{tab:principal} lists shortcuts to compose complete symbols
% of life tables and benefits of insurance and annuities. For space
% considerations only the mandatory arguments are given, but it should
% be noted that all commands of \autoref*{tab:principal} accept the
% optional arguments \meta{ll}, \meta{ul} and \meta{ur} of
% \cmd{\actsymb}.
%
% \begin{table}[t!]
%   \centering
%   \caption{Shortcuts for life table, insurance and annuity
%     symbols. All commands accept arguments the optional arguments
%     \meta{ll}, \meta{ul} and \meta{ur} of \cmd{\actsymb} of
%     \cmd{\actsymb}.}
%   \label{tab:principal}
%   \begin{tabular}{lll}
%       \toprule
%       Definition & Example & Output \\
%       \midrule
%       \cmd{\lx}\marg{age} & \rshowcase{\lx{x}} \\
%       \cmd{\Lx}\marg{age} & \rshowcase{\Lx{x}}$^\dagger$ \\
%       \cmd{\dx}\marg{age} & \rshowcase{\dx[n]{x}} \\
%       \cmd{\Dx}\marg{age} & \rshowcase{\Dx[n]{x}}$^\dagger$ \\
%       \cmd{\px}\marg{age} & \rshowcase{\px[t]{x}} \\
%       \cmd{\qx}\marg{age} & \rshowcase{\qx[t]{x}} \\
%       \cmd{\eringx}\marg{lr} & \rshowcase{\eringx{x:\angln}} \\
%       \addlinespace
%       \cmd{\Ax}\marg{lr}   & \rshowcase{\Ax{x:\angln}} \\
%       \cmd{\Ax*}\marg{lr}  & \rshowcase{\Ax*{x:\angln}} \\
%       \cmd{\Ex}\marg{lr}   & \rshowcase{\Ex[n]{x}} \\
%       \addlinespace
%       \cmd{\ax}\marg{lr}   & \rshowcase{\ax{x:\angln}} \\
%       \cmd{\ax*}\marg{lr}  & \rshowcase{\ax*{x:\angln}} \\
%       \cmd{\ax**}\marg{lr} & \rshowcase{\ax**{x:\angln}} \\
%       \cmd{\sx}\marg{lr}   & \rshowcase{\sx{x:\angln}} \\
%       \cmd{\sx*}\marg{lr}  & \rshowcase{\sx*{x:\angln}} \\
%       \cmd{\sx**}\marg{lr} & \rshowcase{\sx**{x:\angln}} \\
%       \cmd{\aringx}\marg{lr} & \rshowcase{\aringx{x:\angln}} \\
%       \bottomrule
%   \end{tabular} \\
%     \raggedright
%     \hspace*{1em}{\footnotesize $^\dagger$ Calligraphic letters
%     obtained using \cmd{\mathcal}. These specific versions of
%     $\mathcal{L}$ and $\mathcal{D}$ require package \pkg{rsfso}
%     \citep{rsfso}. Standard versions are
%     $\text{\usefont{OMS}{zplm}{m}{n} L}$ and
%     $\text{\usefont{OMS}{zplm}{m}{n} D}$.}
% \end{table}
%
% One will note that shortcuts for insurance benefits come in two
% variants: one for the benefit payable at the end of the year of
% death (standard shortcut) and one for the benefit payable at the
% time of death (|*| variant). Shortcuts for annuity benefits have
% three variants: payable at the end of the period (standard),
% continuously (|*| variant) and at the beginning of the period (|**|
% variant).
%
% \begin{DescribeMacro}{\@actinssc}
% \begin{DescribeMacro}{\@actannsc}
% \begin{DescribeMacro}{\@actothersc}
%   One may define additional shortcut macros similar to those of
%   \autoref{tab:principal} using the internal commands
%   \begin{quote}
%     |\@actinssc|\marg{symbol} \\
%     |\@actannsc|\marg{symbol} \\
%     |\@actothersc|\marg{symbol} \\
%   \end{quote}
%   Command |\@actinssc| automatically provides the two shortcuts for
%   insurance benefits, command |\@actannsc| provides the three
%   shortcuts for annuity benefits and command |\@actothersc|
%   provides a single shortcut for \meta{symbol}.
%
%   For example, the package defines the shortcut \cmd{\px} with the
%   equivalent of %
%   ^^A below we abuse the texample environment to avoid defining
%   ^^A another environment just for this purpose
%   \begin{demo}
%     \begin{texample}[1]
%       |\newcommand{\px}{\@actothersc{p}}| \producing
%     \end{texample}
%   \end{demo}
%   the shortcuts \cmd{\Ax} and \cmd{\Ax*} with
%   \begin{demo}
%     \begin{texample}[1]
%       |\newcommand{\Ax}{\@actinssc{A}}| \producing
%     \end{texample}
%   \end{demo}
%   and shortcuts \cmd{\ax}, \cmd{\ax*} and \cmd{\ax**} with
%   \begin{demo}
%     \begin{texample}[1]
%       |\newcommand{\ax}{\@actannsc{a}}| \producing
%     \end{texample}
%   \end{demo}
% \end{DescribeMacro}
% \end{DescribeMacro}
% \end{DescribeMacro}
%
% \begin{rem}
%   When used in the preamble of a document, the above commands need
%   to be protected by a |\makeatletter| \dots\ |\makeatother| pair.
%   Therefore, to define shortcuts, say, \cmd{\Bx} and \cmd{\Bx*} to
%   the imaginary benefits $B$ and $\bar{B}$, one would write in the
%   preamble
%   \begin{demo}
%     \begin{texample}[1]
%       |\makatletter| \\
%       |\newcommand{\Bx}{\@actinssc{B}}| \\
%       |\makeatother| \producing
%     \end{texample}
%   \end{demo}
% \end{rem}
%
% \subsection{Symbols for premiums, reserves and paid-up insurance}
% \label{sec:shortcuts:premres}
%
% Symbols for premiums, reserves and paid-up insurance come in two
% main variants: without an explicit benefit between parentheses for
% insurances payable at the end of the year of death, and with the
% type of benefit specified for all other insurances and annuities.
%
% \autoref{tab:premres} lists two sets of shortcuts to compose symbols
% for premiums, reserves and paid-up insurance. Shortcut from the
% first set take in mandatory argument the content of the \meta{lr}
% subscript of actuarial symbols. Those from the second set require a
% complete benefit symbol in argument. One may use the shortcuts of
% \autoref{tab:principal} to this effect. Again, only the mandatory
% arguments are given in the table, but all commands accept the
% optional arguments \meta{ll}, \meta{ul} and \meta{ur} of
% \cmd{\actsymb}.
%
% \begin{table}[t]
%   \centering
%   \caption{Shortcuts premiums, reserves and paid-up insurance. All
%     commands accept arguments the optional arguments \meta{ll},
%     \meta{ul} and \meta{ur} of \cmd{\actsymb}.}
%   \label{tab:premres}
%   \begin{tabular*}{1.0\linewidth}{lll}
%       \toprule
%       Definition & Example & Output \\
%       \midrule
%       \cmd{\Px}\marg{lr}  & \rshowcase{\Px[h]{x}} \\
%       \cmd{\Vx}\marg{lr}  & \rshowcase{\Vx[k]{x}} \\
%       \cmd{\Wx}\marg{lr}  & \rshowcase{\Wx[k]{x}} \\
%       \addlinespace
%       \cmd{\premium}\marg{benefit}  & \rshowcase{\premium[t]{\Ax{x}}} \\
%       \cmd{\premium*}\marg{benefit} & \rshowcase{\premium*[t]{\Ax*{x:\angln}}} \\
%       \cmd{\reserve}\marg{benefit}  & \rshowcase{\reserve[t]{\ax**{x}}} \\
%       \cmd{\reserve*}\marg{benefit} & \rshowcase{\reserve*[t]{\ax*{x}}} \\
%       \cmd{\paidup}\marg{benefit}   & \rshowcase{\paidup[k]{\Ax*{x}}} \\
%       \cmd{\paidup*}\marg{benefit}  & \rshowcase{\paidup*[k][h]{\Ax*{x}}} \\
%       \bottomrule
%   \end{tabular*}
% \end{table}
%
% \begin{DescribeMacro}{\@actpremres}
%   The package eases definition of additional shortcut macros similar
%   to the second set of \autoref{tab:premres} using the internal
%   command
%   \begin{quote}
%     |\@actpremres|\marg{principle}
%   \end{quote}
%   It defines shortcuts for discrete (default) or continuous (|*|
%   variant) premiums, reserves and paid-up insurance. For example,
%   the package defines the shortcut the shortcuts \cmd{\premium} and
%   \cmd{\premium*} with
%   \begin{demo}
%     \begin{texample}[1]
%       |\newcommand{\premium}{\@actpremres{P}}| \producing
%     \end{texample}
%   \end{demo}
%   The remark at the end of \autoref{sec:shortcuts:principal} also
%   applies here.
% \end{DescribeMacro}
%
% \subsection{Auxiliary symbols}
% \label{sec:shortcuts:aux}
%
% \autoref{tab:aux} lists shortcuts for a few common auxiliary symbols
% used in the lower right subscript \fbox{\small III}.
%
% \begin{table}
%   \centering
%   \caption{Shortcuts for auxiliary symbols}
%   \label{tab:aux}
%   \begin{tabular*}{1.0\linewidth}{lll}
%       \toprule
%       Definition & Example & Output \\
%       \midrule
%       \cmd{\term}\marg{age}\marg{duration}      & \rshowcase{\Ax{\term{x}{n}}} \\
%       \cmd{\termxn}                             & \rshowcase{\Ax{\termxn}} \\
%       \cmd{\pureendow}\marg{age}\marg{duration} & \rshowcase{\Ax{\pureendow{x}{n}}} \\
%       \cmd{\pureendowxn}                        & \rshowcase{\Ax{\pureendowxn}} \\
%       \cmd{\endow}\marg{age}\marg{duration}     & \rshowcase{\ax**{\endow{x}{n}}} \\
%       \cmd{\endowxn}                            & \rshowcase{\ax**{\endowxn}} \\
%       \cmd{\joint}\marg{ages}$^\dagger$         & \rshowcase{\Ax{\joint{xy}z}} \\
%       \bottomrule
%   \end{tabular*} \\
%   \raggedright
%   \hspace*{1em}{\footnotesize $^\dagger$ alias for \cmd{\overline}}
% \end{table}
%
% \subsection{Symbols for varying benefit insurance and annuities}
% \label{sec:shortcuts:twoletsymb}
%
% \autoref{tab:twoletsymb} lists shortcuts for common two-letter
% symbols of varying benefit insurance and annuities. These shortcuts
% can be used as principal symbol in \cmd{\actsymb} or in the commands
% of \autoref{tab:principal}.
%
% \begin{table}
%   \centering
%   \caption{Shortcuts for varying benefit insurance and annuities}
%   \label{tab:twoletsymb}
%   \begin{tabular}{ll@{\hspace{15mm}}ll}
%       \toprule
%       Definition & Output & Definition & Output \\
%       \midrule
%       \cmd{\IA}   & $\IA$   & \cmd{\DA}   & $\DA$ \\
%       \cmd{\IA*}  & $\IA*$  & \cmd{\DA*}  & $\DA*$ \\
%       \addlinespace[6pt]
%       \cmd{\IbA}   & $\IbA$   & \cmd{\DbA}   & $\DbA$ \\
%       \cmd{\IbA*}  & $\IbA*$  & \cmd{\DbA*}  & $\DbA*$ \\
%       \addlinespace[6pt]
%       \cmd{\ImA}   & $\ImA$   & \cmd{\DmA}   & $\DmA$ \\
%       \cmd{\ImA*}  & $\ImA*$  & \cmd{\DmA*}  & $\DmA*$ \\
%       \addlinespace[6pt]
%       \cmd{\Ia}   & $\Ia$   & \cmd{\Da}   & $\Da$ \\
%       \cmd{\Ia*}  & $\Ia*$  & \cmd{\Da*}  & $\Da*$ \\
%       \cmd{\Ia**} & $\Ia**$ & \cmd{\Da**} & $\Da**$ \\
%       \addlinespace[6pt]
%       \cmd{\Is}   & $\Is$   & \cmd{\Ds}   & $\Ds$ \\
%       \cmd{\Is*}  & $\Is*$  & \cmd{\Ds*}  & $\Ds*$ \\
%       \cmd{\Is**} & $\Is**$ & \cmd{\Ds**} & $\Ds**$ \\
%       \bottomrule
%   \end{tabular}
% \end{table}
%
% \begin{DescribeMacro}{\@twoletinssc}
% \begin{DescribeMacro}{\@twoletannsc}
%   One may define additional shortcuts for two-letter insurance and
%   annuity benefit symbols using the internal commands
%   \begin{quote}
%     |\@twoletinssc|\oarg{length}\marg{symbol\_1}\marg{symbol\_2} \\
%     |\@twoletannsc|\oarg{length}\marg{symbol\_1}\marg{symbol\_2}
%   \end{quote}
%   These commands are similar to those of
%   \autoref{sec:shortcuts:principal}: |\@twoletinssc| provides two
%   shortcuts (standard and |*| variant), whereas |\@twoletannsc|
%   provides three (standard, |*| and |**| variants).
% \end{DescribeMacro}
% \end{DescribeMacro}
%
% \subsection{First, second and third precedence}
% \label{sec:shortcuts:preced}
%
% \autoref{tab:preced} lists shortcuts and aliases for the first,
% second and third precedence numbers, top and bottom. These shortcuts
% can be used in auxiliary symbols in \cmd{\actsymb} or in the
% commands of \autoref{tab:principal}.
%
% \begin{table}[t]
%   \centering
%   \caption{Shortcuts and aliases for precedence numbers}
%   \label{tab:preced}
%   \begin{tabular}{lll}
%       \toprule
%       Definition & Alias & Output \\
%       \midrule
%       \cmd{\firsttop}\marg{status}  & \cmd{\itop}\marg{status}   & $\Ax{\itop{x}:\angln}$ \\
%       \cmd{\secondtop}\marg{status} & \cmd{\iitop}\marg{status}  & $\Ax{x\iitop{y}z}$ \\
%       \cmd{\thirdtop}\marg{status}  & \cmd{\iiitop}\marg{status} & $\Ax{xy\iiitop{z}}$ \\
%       \addlinespace[6pt]
%       \cmd{\firsttop*}\marg{status}  & \cmd{\itop*}\marg{status}   & $\Ax{\itop*{x}:\angln}$ \\
%       \cmd{\secondtop*}\marg{status} & \cmd{\iitop*}\marg{status}  & $\Ax{x\iitop*{y}z}$ \\
%       \cmd{\thirdtop*}\marg{status}  & \cmd{\iiitop*}\marg{status} & $\Ax{xy\iiitop*{z}}$ \\
%       \addlinespace[6pt]
%       \cmd{\firstbottom}\marg{status}  & \cmd{\ibottom}\marg{status}   & $\Ax{\ibottom{x}yz}$ \\
%       \addlinespace[3pt]
%       \cmd{\secondbottom}\marg{status} & \cmd{\iibottom}\marg{status}  & $\Ax{\iibottom{x}yz}$ \\
%       \addlinespace[3pt]
%       \cmd{\thirdbottom}\marg{status}  & \cmd{\iiibottom}\marg{status} & $\Ax{xy\iiibottom{z}}$ \\
%       \addlinespace[6pt]
%       \cmd{\firstbottom*}\marg{status}  & \cmd{\ibottom*}\marg{status}   & $\Ax{\ibottom*{x}yz}$ \\
%       \addlinespace[3pt]
%       \cmd{\secondbottom*}\marg{status} & \cmd{\iibottom*}\marg{status}  & $\Ax{\iibottom*{x}yz}$ \\
%       \addlinespace[3pt]
%       \cmd{\thirdbottom*}\marg{status}  & \cmd{\iiibottom*}\marg{status} & $\Ax{xy\iiibottom*{z}}$ \\
%       \addlinespace[6pt]
%       \bottomrule
%   \end{tabular}
% \end{table}
%
% \section{Package options}
% \label{sec:options}
%
% The package offers two options to control the placement of
% precedence numbers:
% \begin{description}
% \item[\mdseries\ttfamily alignpreced] always align top and bottom
%   precedence numbers; makes \cmd{\nthtop} and \cmd{\nthbottom}
%   equivalent to \cmd{\nthtop*} and \cmd{\nthbottom*},
%   respectively;
% \item[\mdseries\ttfamily compactpreced] always put precedence
%   numbers close to the corresponding status; makes \cmd{\nthtop*}
%   and \cmd{\nthbottom*} equivalent to \cmd{\nthtop} and
%   \cmd{\nthbottom}, respectively.
% \end{description}
% Option |alignpreced| can be useful to simplify entry of uniformly
% positioned precedence numbers. On the other hand, option
% |compactpreced| is merely included as a complement to |alignpreced|.
% We do not recommended --- nor see a reason --- to use it as it
% basically renders inoperative commands \cmd{\nthtop*} and
% \cmd{\nthbottom*}.
%
% \appendix
%
% \newpage
% \section{Comprehensive life contingencies symbol list}
% \label{sec:comprehensive}
%
% This appendix was inspired by \citet{Trivedi:lifecon:2004}. The
% title is a nod to the immensely useful \emph{Comprehensive
% {\LaTeX} Symbol List} of \citet{comprehensive}.
%
% We provide below a table of all the life contingencies symbols used
% in \citet[Appendix~3]{Bowers:2e:1997} along with the code to compose
% them with \pkg{actuarialsymbol} loaded. Often times there are
% various ways to obtain a symbol; we tried to always present the most
% compact and legible option. Not all symbols require features of
% \pkg{actuarialsymbol}.
%
% Calligraphic letters are obtained using \cmd{\mathcal}. The specific
% versions of $\mathcal{L}$ and $\mathcal{D}$ that appear in this
% table require package \pkg{rsfso} \citep{rsfso}, which is
% \emph{not} loaded by \pkg{actuarialsymbol}. Standard versions of
% the calligraphic letters are $\text{\usefont{OMS}{zplm}{m}{n} L}$
% and $\text{\usefont{OMS}{zplm}{m}{n} D}$. See also Table~299 of
% \citet{comprehensive}.
%
% We recommend to define shortcuts for entries marked with $\dagger$ if
% the symbols are to be used repeatedly in a document.
%
% %
% \begingroup
% \catcode`\|=12 ^^A to use | in math, disable as verbatim delimiter
% \begin{longtable}{ll}
%   \toprule
%   Symbol & Code to compose the symbol \\
%   \midrule
%   \endfirsthead
%   \toprule
%   Symbol & Code to compose the symbol \\
%   \midrule
%   \endhead
%   \midrule
%   \endfoot
%   \bottomrule
%   \endlastfoot
%   \showcase{a} \\
%   \showcase{a(x)} \\
%   \showcase{a_x} \\
%   \showcase{\ax{\angl{K}}} \\
%   \showcase{\ax*{\angln}} \\
%   \showcase{\ax*{P_t}} \\
%   \showcase{\ax*{\angl{T}}} \\
%   \showcase{\ax*{x}} \\
%   \showcase{\ax*{r}[h]} \\
%   \showcase{\ax*{x + t}[i]} \\
%   \showcase{\ax*{x + t}[r]} \\
%   \showcase{\ax**{\angl{K + 1}}} \\
%   \showcase{\ax**{x}[(m)]} \\
%   \showcase{\aringx{x}[(m)]} \\
%   \showcase{\ax**{x}[\{m\}]} \\
%   \showcase{\ax**[j]{x}} \\
%   \showcase{\ax**[*]{x}} \\
%   \showcase{\ax{\endowxn}} \\
%   \showcase{\ax*{\endowxn}} \\
%   \showcase{\ax**{\endowxn}} \\
%   \showcase{\ax**{\endowxn}[(m)]} \\
%   \showcase{\aringx{\endowxn}[(m)]} \\
%   \showcase{\ax**{\endowxn}[\{m\}]} \\
%   \showcase{\ax*{\joint\endowxn}} \\
%   \showcase{\ax*[][2]{\endowxn}} \\
%   \showcase{\ax[n|]{x}} \\
%   \showcase{\ax*[n|]{x}} \\
%   \showcase{\ax**[n|]{x}} \\
%   \showcase{\ax**[n|]{x}[(m)]} \\
%   \showcase{\ax*{xy|z}[1]} \\
%   \showcase{\ax**{xy}[(m)]} \\
%   \showcase{\ax**{xy:\angln}} \\
%   \showcase{\ax**[][2]{xy:\angln}} \\
%   \showcase{\ax*{x|y}} \\
%   \showcase{\ax*{\joint{x_1 x_2 x_3}}} \\
%   \showcase{\twoletsymb[0.6pt]{a}{A}(x)}$\;^\dagger$ \\
%   \showcase{\twoletsymb[0.6pt]{a}{A}_t}$\;^\dagger$ \\
%   \showcase{\twoletsymb[0.3pt]{a}{C}_t}$\;^\dagger$ \\
%   \showcase{\twoletsymb[0.6pt]{a}{F}_t}$\;^\dagger$ \\
%   \showcase{\twoletsymb[0.6pt]{a}{U}_t}$\;^\dagger$ \\
%   \showcase{\twoletsymb[0.6pt]{a}{V}(x)}$\;^\dagger$ \\
%   \showcase{\twoletsymb[0.6pt]{a}{V}_t}$\;^\dagger$ \\
%   \showcase{A(h)} \\
%   \showcase{A_t} \\
%   \showcase{\Ax{x}} \\
%   \showcase{\Ax*{x}} \\
%   \showcase{\Ax{x}[(m)]} \\
%   \showcase{\Ax*{x}[\text{PR}]} \\
%   \showcase{\Ax{\termxn}} \\
%   \showcase{\Ax{\endowxn}} \\
%   \showcase{\Ax*{\endowxn}} \\
%   \showcase{\Ax{\pureendowxn}} \\
%   \showcase{\Ax[j]{x}} \\
%   \showcase{\Ax[*]{x}} \\
%   \showcase{\Ax*{\termxn}} \\
%   \showcase{\actsymb{\tilde{A}}{\termxn}}$\;^\dagger$ \\
%   \showcase{\Ax[][2]{\pureendowxn}} \\
%   \showcase{\Ax[][2]{\termxn}} \\
%   \showcase{\Ax*[m|]{x}} \\
%   \showcase{\Ax*[m|n]{x}} \\
%   \showcase{\Ax{xy}} \\
%   \showcase{\Ax{\joint{xy}}} \\
%   \showcase{\Ax{xy}[(m)]} \\
%   \showcase{\Ax*{x\iitop{y}}} \\
%   \showcase{\Ax*{\itop{x}y}} \\
%   \showcase{\Ax{xy:\angln}} \\
%   \showcase{\Ax*{\itop{\joint{xy}}:\angln}} \\
%   \showcase{\Ax[][2]{xy:\angln}} \\
%   \showcase{\Ax*{wx\iitop{y}}} \\
%   \showcase{\Ax*{\joint{x_1 x_2 x_3}}} \\
%   \showcase{\actsymb[k]{\mathit{AS}}{}} \\
%   \showcase{\actsymb[k]{\widehat{\mathit{AS}}}{}} \\
%   \showcase{\actsymb[k]{\widehat{AS}}{}} \\
%   \showcase{\twoletsymb[0pt]{A}{S}_{x + h}}$\;^\dagger$ \\
%   \showcase{(\mathit{AAI})}$\;^\dagger$ \\
%   \addlinespace
%   \showcase{\beta} \\
%   \showcase{\bar{\beta}} \\
%   \showcase{\Gamma(\alpha)} \\
%   \showcase{b(u)} \\
%   \showcase{b_j} \\
%   \showcase{b_t} \\
%   \showcase{b_f(t)} \\
%   \showcase{B_t} \\
%   \showcase{\hat{B}_{x + k}} \\
%   \showcase{\hat{B}_{x + t}^{(3)}} \\
%   \showcase{\hat{B}_{x + t}^{(j)}} \\
%   \showcase{\actsymb[h]{\mathit{BP}}{}}$\;^\dagger$ \\
%   \addlinespace
%   \showcase{c} \\
%   \showcase{c_k} \\
%   \showcase{\hat{c}_k} \\
%   \showcase{c(t)} \\
%   \showcase{C_1} \\
%   \showcase{C_2} \\
%   \showcase{C_3} \\
%   \showcase{C_h} \\
%   \showcase{\actsymb[k]{\mathit{CV}}{}}$\;^\dagger$ \\
%   \addlinespace
%   \showcase{\dx{x}[(j)]} \\
%   \showcase{\dx[n]{x}} \\
%   \showcase{\dx[n]{x}[(j)]} \\
%   \showcase{\dx[n]{x}[(\tau)]} \\
%   \showcase{\actsymb[t]{D}{j}} \\
%   \showcase{\actsymb[k + 1]{D}{}} \\
%   \showcase{\DA_{\termxn}} \\
%   \showcase{\DA*_{\termxn}} \\
%   \showcase{\Dx[n]{x}} \\
%   \showcase{\Dx[n]{x}[(j)]} \\
%   \showcase{\Dx[n]{x}[(\tau)]} \\
%   \addlinespace
%   \showcase{e} \\
%   \showcase{e_{h - 1}} \\
%   \showcase{e_x} \\
%   \showcase{\eringx{x}} \\
%   \showcase{\hat{e}_k} \\
%   \showcase{\eringx{\endowxn}} \\
%   \showcase{e_{xy}} \\
%   \showcase{e_{\joint{xy}}} \\
%   \showcase{\mathrm{E}} \\
%   \showcase{E} \\
%   \showcase{E_0} \\
%   \showcase{\Ex[n]{x}} \\
%   \showcase{\twoletsymb[0.6pt]{E}{S}_{x + h + t}}$\;^\dagger$ \\
%   \showcase{\mathit{ELRA}}$\;^\dagger$ \\
%   \addlinespace
%   \showcase{f} \\
%   \showcase{f(u, t)} \\
%   \showcase{f_S(s)} \\
%   \showcase{F_X(x)} \\
%   \showcase{F_t} \\
%   \showcase{F^{(k)}} \\
%   \showcase{F_S(s)} \\
%   \showcase{\actsymb[k]{F}{}} \\
%   \addlinespace
%   \showcase{G} \\
%   \showcase{\hat{G}} \\
%   \showcase{G(b)} \\
%   \showcase{G(b)} \\
%   \showcase{G(x; \alpha, \beta)} \\
%   \addlinespace
%   \showcase{h(x)} \\
%   \showcase{H(r)} \\
%   \showcase{H(x; \alpha, \beta, x_0)} \\
%   \showcase{\actsymb[u]{\twoletsymb[0.6pt]{h}{p}}{x + t}[(\tau)]}$\;^\dagger$ \\
%   \showcase{\actsymb{\twoletsymb[0.6pt]{h}{\mu}}{x + t}[(j)](u)}$\;^\dagger$ \\
%   \addlinespace
%   \showcase{i'_{k + 1}} \\
%   \showcase{\hat{\imath}_{k + 1}} \\
%   \showcase{i(s, s + t)} \\
%   \showcase{I_k} \\
%   \showcase{I_d} \\
%   \showcase{I_d(x)} \\
%   \showcase{\actsymb[j]{i}{k}} \\
%   \showcase{\IA_x} \\
%   \showcase{\IA*_x} \\
%   \showcase{\IbA*_x} \\
%   \showcase{\ImA*_x} \\
%   \showcase{\IA_{\termxn}} \\
%   \addlinespace
%   \showcase{J} \\
%   \showcase{j(s, s + t, s + u} \\
%   \addlinespace
%   \showcase{\actsymb[t]{\bar{k}}{x}} \\
%   \showcase{K} \\
%   \showcase{K(x)} \\
%   \showcase{K(xy)} \\
%   \showcase{K(\joint{xy})} \\
%   \addlinespace
%   \showcase{\lx{x}} \\
%   \showcase{\lx{[x] + k}} \\
%   \showcase{\lx{x}[(\tau)]} \\
%   \showcase{l(x, u)} \\
%   \showcase{l_f(x, u)} \\
%   \showcase{L} \\
%   \showcase{L_1} \\
%   \showcase{L_x} \\
%   \showcase{L(h)} \\
%   \showcase{\actsymb[t]{L}{}} \\
%   \showcase{\actsymb[t]{L}{}[\underline{2}]} \\
%   \showcase{\actsymb[t]{L}{e}} \\
%   \showcase{\actsymb[t]{L}{e}[\underline{2}]} \\
%   \showcase{\mathcal{L}(x)} \\
%   \showcase{\Lx{x}[(\tau)]} \\
%   \addlinespace
%   \showcase{m(x)} \\
%   \showcase{m_x} \\
%   \showcase{m_x^{(j)}} \\
%   \showcase{m_x^{(\tau)}} \\
%   \showcase{m_x^{\prime (j)}} \\
%   \showcase{\mathrm{M}_X(t)} \\
%   \showcase{M(x)} \\
%   \addlinespace
%   \showcase{n(u)} \\
%   \showcase{N} \\
%   \showcase{N(t)} \\
%   \addlinespace
%   \showcase{p(j)} \\
%   \showcase{p(x)} \\
%   \showcase{p_k} \\
%   \showcase{\px{[x] + r}} \\
%   \showcase{p^{*n}(x)} \\
%   \showcase{\px[t]{x}} \\
%   \showcase{\px[t]{x}[(\tau)]} \\
%   \showcase{\px[t]{x}[\prime (j)]} \\
%   \showcase{\px[t]{xy}} \\
%   \showcase{\px[t]{\joint{xy}}} \\
%   \showcase{\px[u]{\joint{xy} + t}} \\
%   \showcase{\px[t]{\nthtop{\hphantom{x_1x_2}k}{\joint{x_1x_2x_3}}}} \\
%   \showcase{P(x)} \\
%   \showcase{P(s, t)} \\
%   \showcase{P_t} \\
%   \showcase{\actsymb[][T]{P}{t}} \\
%   \showcase{P^a} \\
%   \showcase{P_x} \\
%   \showcase{\actsymb[j]{P}{x}} \\
%   \showcase{\actsymb[*]{P}{x}} \\
%   \showcase{\Px{\endowxn}[A]} \\
%   \showcase{\Px{\endowxn}} \\
%   \showcase{\Px{\joint{xy}}} \\
%   \showcase{\Px{\termxn}} \\
%   \showcase{\Px{\pureendowxn}} \\
%   \showcase{P^{*n}(x)} \\
%   \showcase{\Px[h]{x}} \\
%   \showcase{\Px[h]{\endowxn}} \\
%   \showcase{\twoletsymb[0.6pt]{P}{a}(x)} \\
%   \showcase{\twoletsymb[0.6pt]{P}{a}_t} \\
%   \showcase{\premium{\Ax*{\endowxn}}} \\
%   \showcase{\premium{\ax**[n|]{x}}} \\
%   \showcase{\premium{\ax*[n|]{x}}} \\
%   \showcase{\premium{\Ax*{x}}[(m)]} \\
%   \showcase{\premium{\Ax*{x}}[\{m\}]} \\
%   \showcase{\premium{\Ax*{x}[\text{PR}]}} \\
%   \showcase{\premium*{\Ax*{\endowxn}}} \\
%   \showcase{\premium*{\Ax*{\termxn}}} \\
%   \showcase{\premium*{\Ax*{\pureendowxn}}} \\
%   \showcase{\premium{\Ax*{\endowxn}}[(m)]} \\
%   \showcase{\premium{\Ax*{\termxn}}[(m)]} \\
%   \showcase{\premium*[h]{\Ax*{x}}} \\
%   \showcase{\premium*[h]{\Ax*{\endowxn}}} \\
%   \showcase{\premium[h]{\Ax*{x}}[(m)]} \\
%   \showcase{\premium[h]{\Ax*{\endowxn}}[(m)]} \\
%   \showcase{\premium[h]{\Ax*{\endowxn}}[\{m\}]} \\
%   \showcase{\premium{\Ax*{\nthtop{\hphantom{xy}2}{\joint{xyz}}}}} \\
%   \showcase{\premium{\Ax*{\ibottom{x}\iitop{y}z}}} \\
%   \showcase{\tilde{P}_\termxn} \\
%   \addlinespace
%   \showcase{\qx{[x] + r}} \\
%   \showcase{\qx{x}[(d)]} \\
%   \showcase{\qx{x}[(i)]} \\
%   \showcase{\qx{x}[(r)]} \\
%   \showcase{\qx{x}[(w)]} \\
%   \showcase{\actsymb{\hat{q}}{x + k}[(j)]} \\
%   \showcase{\qx{xy}} \\
%   \showcase{\qx[k|]{x}} \\
%   \showcase{\qx[t]{x}} \\
%   \showcase{\qx[t]{x}[(j)]} \\
%   \showcase{\qx[t]{x}[(\tau)]} \\
%   \showcase{\qx[t]{x}[\prime (j)]} \\
%   \showcase{\qx[t|u]{x}} \\
%   \showcase{\qx[n]{\itop{x}y}} \\
%   \showcase{\qx[n]{x\iitop{y}}} \\
%   \showcase{\qx[k|]{xy}} \\
%   \showcase{\qx[n]{\ibottom{x}\iitop{y}z}} \\
%   \showcase{\qx[\infty]{\ibottom*{x}\iibottom*{y}\iiitop{z}}} \\
%   \addlinespace
%   \showcase{r} \\
%   \showcase{r_C} \\
%   \showcase{r_F} \\
%   \showcase{r_N} \\
%   \showcase{\twoletsymb[0.6pt]{r}{A}_t} \\
%   \showcase{\twoletsymb[0.6pt]{r}{F}_t} \\
%   \showcase{\twoletsymb[0.6pt]{r}{V}_t} \\
%   \showcase{R} \\
%   \showcase{\tilde{R}} \\
%   \showcase{R(x, h, y)} \\
%   \addlinespace
%   \showcase{s(x)} \\
%   \showcase{\sx**{\angln}} \\
%   \showcase{s(x,u)} \\
%   \showcase{\sx*{\endowxn}} \\
%   \showcase{\sx**{\endowxn}} \\
%   \showcase{S} \\
%   \showcase{S(t)} \\
%   \showcase{S_y} \\
%   \showcase{\actsymb[k]{\mathit{SC}}{}} \\
%   \addlinespace
%   \showcase{T} \\
%   \showcase{\tilde{T}} \\
%   \showcase{T_x} \\
%   \showcase{T_{xy}} \\
%   \showcase{T_{\joint{xy}}} \\
%   \addlinespace
%   \showcase{U(t)} \\
%   \showcase{U_t} \\
%   \showcase{\hat{U}{t}} \\
%   \addlinespace
%   \showcase{\tilde{V}_t} \\
%   \showcase{\Vx{t}} \\
%   \showcase{\Vx[k]{t}} \\
%   \showcase{\Vx[k]{\endowxn}} \\
%   \showcase{\Vx[k]{\termxn}} \\
%   \showcase{\Vx[k]{\pureendowxn}} \\
%   \showcase{\Vx[k]{x}[\text{FPT}]} \\
%   \showcase{\Vx*{t}} \\
%   \showcase{\actsymb[k]{\bar{V}}{t}} \\
%   \showcase{\actsymb[k]{\bar{V}}{\endowxn}} \\
%   \showcase{\actsymb[k]{\bar{V}}{\termxn}} \\
%   \showcase{\actsymb[k]{\bar{V}}{\pureendowxn}} \\
%   \showcase{\actsymb[k]{\bar{V}}{x}[FPT]} \\
%   \showcase{\actsymb[k][h]{V}{\termxn}} \\
%   \showcase{\actsymb[k][h]{V}{\endowxn}[(m)]} \\
%   \showcase{\reserve[k]{\ddot{a}}{x}} \\
%   \showcase{\reserve*[k]{\bar{a}}{x}} \\
%   \showcase{\reserve[k]{\bar{A}}{x}} \\
%   \showcase{\reserve*[k]{\bar{A}}{x}} \\
%   \showcase{\reserve[k]{\bar{A}}{\endowxn}} \\
%   \showcase{\reserve*[k]{\bar{A}}{\endowxn}} \\
%   \showcase{\reserve[k]{\bar{A}}{\termxn}} \\
%   \showcase{\reserve[k]{\bar{A}}{\pureendowxn}} \\
%   \showcase{\reserve[k]{\bar{A}}{\joint{xy}}} \\
%   \showcase{\reserve[k]{\bar{A}}{x}[PR]} \\
%   \showcase{\reserve[k][h]{\bar{A}}{\termxn}} \\
%   \showcase{\reserve[k][h]{\bar{A}}{x}} \\
%   \showcase{\reserve*[k][h]{\bar{A}}{\endowxn}} \\
%   \showcase{\reserve[k][h]{\bar{A}}{\endowxn}[(m)]} \\
%   \addlinespace
%   \showcase{\actsymb{W}{t}} \\
%   \showcase{\actsymb[k]{W}{}} \\
%   \showcase{\actsymb[k]{W}{x}} \\
%   \showcase{\actsymb[k]{W}{\endowxn}} \\
%   \showcase{\actsymb[k][h]{W}{x}} \\
%   \showcase{\twoletsymb[0.6pt]{W}{a}_t} \\
%   \showcase{\paidup*[k]{\bar{A}}{x}} \\
%   \showcase{\paidup*[k]{\bar{A}}{\endowxn}} \\
%   \showcase{\paidup*[k][h]{\bar{A}}{x}} \\
%   \addlinespace
%   \showcase{\actsymb[m]{Z}{t}}\\
%   \addlinespace
%   \showcase{\delta} \\
%   \showcase{\actsymb{\delta}{t}}\\
%   \showcase{\actsymb{\mu}{x}} \\
%   \showcase{\actsymb{\mu}{x}(t)} \\
%   \showcase{\actsymb{\mu}{x}[(i)]}\\
%   \showcase{\actsymb{\mu}{x}[(j)](t)} \\
%   \showcase{\actsymb{\mu}{xy} (t)}\\
%   \showcase{\actsymb{\mu}{\joint{xy}}(t)} \\
%   \showcase{\omega}\\
% \end{longtable}
% \endgroup
%
% \StopEventually{
% \newpage
% \begin{thebibliography}{9}
% \bibitem[{A}merican~{M}athematical {S}ociety(2002)]{amsmath}
%   {A}merican~{M}athematical {S}ociety.
%   \newblock \emph{User’s Guide for the amsmath Package}, 2002.
%   \newblock URL \url{http://www.ctan.org/pkg/amsmath/}.
%
% \bibitem[Bowers et~al.(1997)Bowers, Gerber, Hickman, Jones, and
%   Nesbitt]{Bowers:2e:1997}
%   N.~L. Bowers, H.~U. Gerber, J.~C. Hickman, D.~A. Jones, and C.~J. Nesbitt.
%   \newblock \emph{Actuarial Mathematics}.
%   \newblock Society of Actuaries, Shaumburg, IL, second edition, 1997.
%   \newblock ISBN 0-9389594-6-8.
%
% \bibitem[Goulet(2017)]{actuarialangle}
%   V. Goulet.
%   \newblock \emph{Actuarial angle symbol for life contingencies and financial
%     mathematics}, 2017.
%   \newblock URL \url{http://www.ctan.org/pkg/actuarialangle/}.
%
% \bibitem[Høgholm and Madsen(2015)]{mathtools}
%   M. Høgholm and L. Madsen.
%   \newblock \emph{The mathtools package}, 2015.
%   \newblock URL \url{http://www.ctan.org/pkg/mathtools/}.
%
% \bibitem[Pakin(2015)]{comprehensive}
%   S. Pakin.
%   \newblock \emph{The Comprehensive {\LaTeX} Symbol List}, 2015.
%   \newblock URL \url{http://www.ctan.org/pkg/comprehensive/}.
%
% \bibitem[Robertson(2007)]{subdepth}
%   W. Robertson.
%   \newblock \emph{Unify subscript depths}, 2007.
%   \newblock URL \url{http://www.ctan.org/pkg/subdepth/}.
%
% \bibitem[Sharpe(2015)]{rsfso}
%   M. Sharpe.
%   \newblock \emph{An acute script font based on {rsfs}}, 2015.
%   \newblock URL \url{http://www.ctan.org/pkg/rsfso/}.
%
% \bibitem[Trivedi(2004)]{Trivedi:lifecon:2004}
%   E. Trivedi.
%   \newblock \emph{Life Contingencies' Symbols}, 2004.
%   \newblock \texttt{lifecon} 2.1 User Guide.
%
% \bibitem[Wolthuis(2004)]{Wolthuis:notation:2004}
%   H. Wolthuis.
%   \newblock International actuarial notation.
%   \newblock In Jozef Teugels and Bj{\o}rn Sundt, editors, \emph{Encyclopedia of
%   Actuarial Science}. Wiley, 2004.
%   \newblock ISBN 978-0-47001250-5.
%   \newblock URL \url{http://onlinelibrary.wiley.com/book/10.1002/9780470012505}.
% \end{thebibliography}
% }
%
% \newpage
% \section{Implementation}
% \label{sec:implementation}
%
% This appendix contains the annotated source code of the package.
% Most readers can stop reading here.
%\iffalse
%<*package>
%\fi
%
% \subsection{Required packages}
% \label{sec:implementation:packages}
%
% \begin{macro}{RequirePackage}
%   The package depends on packages \pkg{amsmath} \citep{amsmath} and
%   \pkg{actuarialangle} \citep{actuarialangle}.
%    \begin{macrocode}
\RequirePackage{amsmath}
\RequirePackage{actuarialangle}
%    \end{macrocode}
% \end{macro}
%
% \subsection{Package options}
% \label{sec:implementation:options}
%
% \begin{macro}{\ifacts@alignpreced}
% \begin{macro}{\ifacts@acompactpreced}
%   We define two flags to keep track of the type of alignment for
%   precedence numbers. Both are false by default.
%    \begin{macrocode}
\newif\ifacts@alignpreced   \acts@alignprecedfalse
\newif\ifacts@compactpreced \acts@compactprecedfalse
%    \end{macrocode}
% \end{macro}
% \end{macro}
% \begin{macro}{\DeclareOption}
%   Declaration of the package options and processing.
%    \begin{macrocode}
\DeclareOption{alignpreced}{\acts@alignprecedtrue}
\DeclareOption{compactpreced}{\acts@compactprecedtrue}
\ProcessOptions
%    \end{macrocode}
% \end{macro}
%
% \subsection{Actuarial symbol}
% \label{sec:implementation:actsymb}
%
% We face three main challenges when trying to compose actuarial
% symbols: position a subscript and a superscipt to the left of a main
% symbol; ensure that the left and right subscripts and superscripts
% are vertically aligned; maintain consistency with the placement of
% subscripts and superscripts elsewhere in the document.
%
% \subsubsection{User level and parsing commands}
% \label{sec:implementation:actsymb:user}
%
% \begin{macro}{actsymb}
%   We first define the user level command.
%    \begin{macrocode}
\DeclareRobustCommand{\actsymb}{\acts@actsymb}
%    \end{macrocode}
% \end{macro}
% \begin{macro}{\acts@actsymb}
% \begin{macro}{\acts@@actsymb}
% \begin{macro}{\acts@@@actsymb}
% \begin{macro}{\acts@@@@actsymb}
%   With its unusual syntax, \cmd{\actsymb} requires four parsing
%   macros to pass the arguments to the actual workhorse
%   \cmd{\@actsymbol}.
%    \begin{macrocode}
\newcommand\acts@actsymb{%
  \@ifnextchar[{\acts@@actsymb}%
               {\acts@@actsymb[]}}
\newcommand\acts@@actsymb{}
\def\acts@@actsymb[#1]{%
  \@ifnextchar[{\acts@@@actsymb[#1]}%
               {\acts@@@actsymb[#1][]}}
\newcommand\acts@@@actsymb{}
\def\acts@@@actsymb[#1][#2]{%
  \@ifnextchar[{\acts@@@@actsymb[#1][#2]}%
               {\acts@@@@actsymb[#1][#2][]}}
\newcommand\acts@@@@actsymb{}
\def\acts@@@@actsymb[#1][#2][#3]#4#5{%
  \@ifnextchar[{\@actsymbol[#1][#2][#3]{#4}{#5}}%
               {\@actsymbol[#1][#2][#3]{#4}{#5}[]}}
%    \end{macrocode}
% \end{macro}
% \end{macro}
% \end{macro}
% \end{macro}
%
% \subsubsection{Symbol construction}
% \label{sec:implementation:actsymb:workhorse}
%
% \begin{macro}{\@actsymbol}
%   We now define the workhorse to compose an actuarial symbol
%   allowing for auxiliary symbols on all four corners around a
%   principal symbol. The macro \cmd{\@actsymbol} has the following
%   six arguments, using the notation from of the schematics
%   representations \eqref{schema1} and \eqref{schema2}:
%   \begin{enumerate}[nosep]
%   \item lower left subscript \fbox{\small I}; may be empty;
%   \item upper left superscript \fbox{\small II}; may be empty;
%   \item principle symbol $P$; when empty a standard symbol is
%     typeset;
%   \item principal symbol $S$; assumed present;
%   \item lower right subscript \fbox{\small III}; assumed present;
%   \item upper right superscript \fbox{\small IV}; may be empty.
%   \end{enumerate}
%   The definition below is heavily inspired by the code of
%   \cmd{\prescript} from package \pkg{mathtools} \citep{mathtools}
%   which, as reported by the author, is itself based on a posting to
%   |comp.text.tex| by Michael J.~Downes. Some of the comments below
%   are still his.
%    \begin{macrocode}
\def\@actsymbol[#1][#2][#3]#4#5[#6]{
%    \end{macrocode}
%   We put the principle symbol and all subscripts and superscript into
%   horizontal boxes using \cmd{\@mathmeasure} of package \pkg{amsmath}.
%    \begin{macrocode}
  \@mathmeasure\z@\displaystyle{#3}
  \@mathmeasure\@ne\scriptstyle{#1}
  \@mathmeasure\tw@\scriptstyle{#2}
  \@mathmeasure\thr@@\scriptstyle{#5}
  \@mathmeasure4\scriptstyle{#6}
%    \end{macrocode}
%   Do not let a preceding mathord symbol approach without any
%   intervening space.
%    \begin{macrocode}
  \mathord{}
%    \end{macrocode}
%   Right align the left subscript and superscript (if there is a
%   superscript) by adjusting the width of their horizontal boxes.
%    \begin{macrocode}
  \ifdim \wd\tw@>\z@ \ifdim \wd\tw@>\wd\@ne
    \setbox\@ne\hbox to\wd\tw@{\hfil\unhbox\@ne}
  \else
    \setbox\tw@\hbox to\wd\@ne{\hfil\unhbox\tw@}
  \fi\fi
%    \end{macrocode}
%   Bottom align the left and right subscripts by making their heights
%   equal.
%    \begin{macrocode}
  \ifdim \ht\@ne>\ht\thr@@ \ht\thr@@=\ht\@ne\else \ht\@ne=\ht\thr@@\fi
%    \end{macrocode}
%   Actual symbol construction. First the case where there is at least
%   one superscript (the nested conditionals emulate an OR condition).
%    \begin{macrocode}
  \ifnum\ifdim \wd\tw@>\z@ 1\else\ifdim \wd4>\z@ 1\else 0\fi\fi
    =1
%    \end{macrocode}
%   To ensure equal spacing with the subscripts on both sides of the
%   principal symbol, make the depths of the left and right
%   supercripts equal.
%    \begin{macrocode}
    \ifdim \dp\tw@>\dp4 \dp4=\dp\tw@\else \dp\tw@=\dp4\fi
%    \end{macrocode}
%   Typesetting of the left subscript and superscript. Use
%   \cmd{\mathopen} to suppress space between those and the principal
%   symbol even when the latter is not of type ord.
%    \begin{macrocode}
    \mathopen{\vphantom{#4}^{\box\tw@}\sb{\box\@ne}}%
%    \end{macrocode}
%   Typesetting of the principal symbol and the right subscript and
%   superscript. The principal symbol is either a principal symbol alone
%   or a principle symbol and a benefit symbol between parentheses.
%    \begin{macrocode}
    \ifdim\wd\z@=\z@
      #4^{\box4}\sb{\box\thr@@}%
    \else
      #3^{\box4}(#4^{}\sb{\box\thr@@})%
    \fi
  \else
%    \end{macrocode}
%   Same ideas as above in the simpler case without superscripts. That
%   way, {\TeX}'s default placement of subscripts and superscripts is
%   maintained.
%    \begin{macrocode}
    \mathopen{\vphantom{#4}\sb{\box\@ne}}%
    \ifdim\wd\z@=\z@
      #4\sb{\box\thr@@}%
    \else
      #3(#4\sb{\box\thr@@})%
    \fi
  \fi
}
%    \end{macrocode}
% \end{macro}
%
% \subsubsection{Shortcuts for basic symbols of life tables, insurance and annuities}
% \label{sec:implementation:actsymb:principalsc}
%
% \begin{macro}{\@actinssc}
% \begin{macro}{\@actannsc}
% \begin{macro}{\@actothersc}
%   We first define the internal macros and parsing macros used to
%   define shortcuts for insurance benefits (\cmd{\@actinssc}),
%   annuity benefits (\cmd{\@actannsc}) and simple symbols
%   (\cmd{\@actothersc}).
%    \begin{macrocode}
\newcommand\@actinssc[1]{%
  \@ifstar{\acts@actsc{\bar{#1}}}%
          {\acts@actsc{#1}}}
\newcommand\@actannsc[1]{%
  \@ifstar{\@ifstar{\acts@actsc{\ddot{#1}}}%
                   {\acts@actsc{\bar{#1}}}}
          {\acts@actsc{#1}}}
\newcommand\acts@actsc[1]{%
  \@ifnextchar[{\acts@@actsc{#1}}%
               {\acts@@actsc{#1}[]}}
\let\@actothersc\acts@actsc
\newcommand\acts@@actsc{}
\def\acts@@actsc#1[#2]{%
  \@ifnextchar[{\acts@@@actsc{#1}[#2]}%
               {\acts@@@actsc{#1}[#2][]}}
\newcommand\acts@@@actsc{}
\def\acts@@@actsc#1[#2][#3]#4{%
  \@ifnextchar[{\@actsymbol[#2][#3][]{#1}{#4}}%
               {\@actsymbol[#2][#3][]{#1}{#4}[]}}
%    \end{macrocode}
% \end{macro}
% \end{macro}
% \end{macro}
% \begin{macro}{\DeclareRobustCommand}
%   We can now define all the shortcuts of \autoref{tab:principal}.
%    \begin{macrocode}
\DeclareRobustCommand{\lx}{\@actothersc {\ell}}
\DeclareRobustCommand{\Lx}{\@actothersc {\mathcal{L}}}
\DeclareRobustCommand{\dx}{\@actothersc d}
\DeclareRobustCommand{\Dx}{\@actothersc{\mathcal{D}}}
\DeclareRobustCommand{\px}{\@actothersc p}
\DeclareRobustCommand{\qx}{\@actothersc q}
\DeclareRobustCommand{\Ax}{\@actinssc A}
\DeclareRobustCommand{\Ex}{\@actinssc E}
\DeclareRobustCommand{\ax}{\@actannsc a}
\DeclareRobustCommand{\sx}{\@actannsc s}
\DeclareRobustCommand{\aringx}{\@actothersc{\mathring{a}}}
\DeclareRobustCommand{\eringx}{\@actothersc{\mathring{e}}}
%    \end{macrocode}
% \end{macro}
%
% \subsubsection{Shortcuts for premiums, reserves and paid-up insurance}
% \label{sec:implementation:actsymb:premressc}
%
% \begin{macro}{\@actpremres}
%   Not unlike above, we first define the internal and parsing macros
%   used to create shortcut for premiums, reserves and paid-up insurance
%   symbols.
%    \begin{macrocode}
\newcommand\@actpremres[1]{%
  \@ifstar{\acts@actpremres{\bar{#1}}}%
          {\acts@actpremres{#1}}}
\newcommand\acts@actpremres[1]{%
  \@ifnextchar[{\acts@@actpremres{#1}}%
               {\acts@@actpremres{#1}[]}}
\newcommand\acts@@actpremres{}
\def\acts@@actpremres#1[#2]{%
  \@ifnextchar[{\acts@@@actpremres{#1}[#2]}%
               {\acts@@@actpremres{#1}[#2][]}}
\newcommand\acts@@@actpremres{}
\def\acts@@@actpremres#1[#2][#3]#4{%
  \@ifnextchar[{\@actsymbol[#2][#3][#1]{{#4}}{}}%
               {\@actsymbol[#2][#3][#1]{{#4}}{}[]}}
%    \end{macrocode}
% \end{macro}
% \begin{macro}{\DeclareRobustCommand}
%   Then we can define the shortcuts of \autoref{tab:premres}.
%    \begin{macrocode}
\DeclareRobustCommand{\Px}{\@actothersc P}
\DeclareRobustCommand{\Vx}{\@actothersc V}
\DeclareRobustCommand{\Wx}{\@actothersc W}
\DeclareRobustCommand{\premium}{\@actpremres P}
\DeclareRobustCommand{\reserve}{\@actpremres V}
\DeclareRobustCommand{\paidup}{\@actpremres W}
%    \end{macrocode}
% \end{macro}
%
% \subsubsection{Shortcuts for auxiliary symbols}
% \label{sec:implementation:actsymb:auxsc}
%
% \begin{macro}{\DeclareRobustCommand}
%   Here we define the shortcuts for common auxiliary symbols listed in
%   \autoref{tab:aux}. Command \cmd{\joint} is simply defined as an
%   alias of \cmd{\overline}.
%    \begin{macrocode}
\DeclareRobustCommand{\term}[2]{\itop{#1}:\angl{#2}}
\DeclareRobustCommand{\termxn}{{\itop{x}:\angln}}
\DeclareRobustCommand{\pureendow}[2]{#1:\itop{\angl{#2}}}
\DeclareRobustCommand{\pureendowxn}{{x:\itop{\angln}}}
\DeclareRobustCommand{\endow}[2]{#1:\angl{#2}}
\DeclareRobustCommand{\endowxn}{{x:\angln}}
\let\joint\overline
%    \end{macrocode}
% \end{macro}
%
% \subsection{Two-letter symbols}
% \label{sec:implementation:twoletsymb}
%
% \begin{macro}{\twoletkern}
% \begin{macro}{\twoletsymb}
%   The definition of the user level command \cmd{\twoletsymb} is
%   pretty straighforward.
%    \begin{macrocode}
\newlength{\twoletkern}
\setlength{\twoletkern}{1.2pt}
\DeclareRobustCommand{\twoletsymb}[3][\twoletkern]{(#2\kern-#1#3)}
%    \end{macrocode}
% \end{macro}
% \end{macro}
% \begin{macro}{\@twoletinssc}
% \begin{macro}{\@twoletannsc}
%   Following the same lines as above, we define two internal macros
%   to create shortcuts for two-letter insurance and annuity benefit
%   symbols.
%    \begin{macrocode}
\newcommand\@twoletinssc[3][\twoletkern]{%
  \@ifstar{\twoletsymb[#1]{#2}{\bar{#3}}}%
          {\twoletsymb[#1]{#2}{#3}}}
\newcommand\@twoletannsc[3][\twoletkern]{%
  \@ifstar{\@ifstar{\twoletsymb[#1]{#2}{\ddot{#3}}}%
                   {\twoletsymb[#1]{#2}{\bar{#3}}}}%
          {\twoletsymb[#1]{#2}{#3}}}
%    \end{macrocode}
% \end{macro}
% \end{macro}
% \begin{macro}{\newcommand}
%   Follow the shortcuts of \autoref{tab:twoletsymb}. Note that we
%   used a smaller negative kerning adjustment for symbols annuity
%   symbols.
%    \begin{macrocode}
\newcommand{\IA}{\@twoletinssc IA}
\newcommand{\IbA}{\@twoletinssc{\bar{I}}A}
\newcommand{\ImA}{\@twoletinssc{I^{(m)}}A}
\newcommand{\DA}{\@twoletinssc DA}
\newcommand{\DbA}{\@twoletinssc{\bar{D}}A}
\newcommand{\DmA}{\@twoletinssc{D^{(m)}}A}
\newcommand{\Ia}{\@twoletannsc[0.5\twoletkern]Ia}
\newcommand{\Is}{\@twoletannsc[0.5\twoletkern]Is}
\newcommand{\Da}{\@twoletannsc[0.5\twoletkern]Da}
\newcommand{\Ds}{\@twoletannsc[0.5\twoletkern]Ds}
%    \end{macrocode}
% \end{macro}
%
% \subsection{Precedence numbers}
% \label{sec:implementation:preced}
%
% The system of precedence numbers builds on a backbone that used to
% be part of package \pkg{actuarialangle}. The author of the code is
% unknown. Some of the comments on \cmd{\@precendence}, below, are his
% or hers.
%
% \subsubsection{User level lengths and commands}
% \label{sec:implementation:preced:user}
%
% \begin{macro}{\nthtopsep}
% \begin{macro}{\nthtopskip}
% \begin{macro}{\nthbottomsep}
% \begin{macro}{\nthbottomskip}
%   Let us start by defining the lengths needed to position the
%   precedence numbers above or below the statuses.
%    \begin{macrocode}
\newlength{\nthtopsep}
\setlength{\nthtopsep}{2\p@}
\newlength{\nthtopskip}
\setlength{\nthtopskip}{7\p@}
\newlength{\nthbottomsep}
\setlength{\nthbottomsep}{2\p@}
\newlength{\nthbottomskip}
\setlength{\nthbottomskip}{9\p@}
%    \end{macrocode}
% \end{macro}
% \end{macro}
% \end{macro}
% \end{macro}
% \begin{macro}{\nthtop}
% \begin{macro}{\nthbottom}
%   Next up are the two user level commands and the internal commands
%   that call the workhorse \cmd{\@precedence}.
%    \begin{macrocode}
\DeclareRobustCommand{\nthtop}{\@ifstar\@varnthtop\@nthtop}
\DeclareRobustCommand{\nthbottom}{\@ifstar\@varnthbottom\@nthbottom}
\newcommand{\@nthtop}[3][\nthtopsep]{%
  \mathpalette{\@precedence t\z@{#1}{#2}}{#3}}
\newcommand{\@varnthtop}[3][\nthtopskip]{%
  \mathpalette{\@precedence t{#1}\z@{#2}}{#3}}
\newcommand{\@nthbottom}[3][\nthbottomsep]{%
  \mathpalette{\@precedence b\z@{#1}{#2}}{#3}}
\newcommand{\@varnthbottom}[3][\nthbottomskip]{%
  \mathpalette{\@precedence b{#1}\z@{#2}}{#3}}
%    \end{macrocode}
% \end{macro}
% \end{macro}
%
% \begin{macro}{\@nthtop}
% \begin{macro}{\@nthbottom}
%   With the internal commands defined we can now process the package
%   options. Option |alignpreced| makes the internal macros
%   \cmd{\@nthtop} and \cmd{\@nthbottom} equivalent to
%   \cmd{\@varnthtop} and \cmd{\@varnthbottom}, respectively.
%    \begin{macrocode}
\ifacts@alignpreced
  \let\@nthtop\@varnthtop
  \let\@nthbottom\@varnthbottom
\fi
%    \end{macrocode}
% \end{macro}
% \end{macro}
% \begin{macro}{\@varnthtop}
% \begin{macro}{\@varnthbottom}
%   Option |compactpreced| does the opposite.
%    \begin{macrocode}
\ifacts@compactpreced
  \let\@varnthtop\nthtop
  \let\@varnthbottom\nthbottom
\fi
%    \end{macrocode}
% \end{macro}
% \end{macro}
%
% \subsubsection{Number and status positioning}
% \label{sec:implementation:preced:workhorse}
%
% \begin{macro}{\@precedence}
%   We now define the workhorse to position precedence numbers above
%   or below statuses, smashed so that the height of the auxiliary
%   symbol containing the call remains unaffected. The macro
%   \cmd{\@precedence} has the following six arguments:
%   \begin{enumerate}[nosep]
%   \item |t| for top or |b| for bottom;
%   \item baseline skip between the top and bottom symbols;
%   \item spacing between the top and bottom symbols;
%   \item first argument from the user;
%   \item one of math style selecting commands (\cmd{\displaystyle},
%     \cmd{\textstyle}, \cmd{\scriptstyle}, \cmd{\scriptscriptstyle})
%     from \cmd{\mathpalette};
%   \item second argument given by the user.
%   \end{enumerate}
%   This peculiar ordering of the arguments is done to work around the
%   restriction of \cmd{\mathpalette} that it only reads two
%   arguments.
%    \begin{macrocode}
\def\@precedence#1#2#3#4#5#6{%
  \setbox\tw@\hbox{$\m@th#5#4$}%
  \setbox\z@\hbox{$\m@th#5#6$}
  \dimen@\wd\z@
  \vbox to\ht\z@{%
    \baselineskip=#2
    \lineskip=#3
    \lineskiplimit\lineskip
    \if b#1\relax\box\z@\else\vss\fi
    \hbox to\dimen@{\hss\unhbox\tw@\hss}%
    \if t#1\relax\box\z@\else\vss\fi
  }
}
%    \end{macrocode}
% \end{macro}
%
% \subsubsection{Shortcuts and aliases}
% \label{sec:implementation:preced:shortcuts}
%
% \begin{macro}{\DeclareRobustCommand}
%   We finish off by defining the shortcuts and aliases of
%   \autoref{tab:preced}.
%    \begin{macrocode}
\DeclareRobustCommand{\firsttop}{%
  \@ifstar{\@varnthtop{1}}{\@nthtop{1}}}
\DeclareRobustCommand{\secondtop}{%
  \@ifstar{\@varnthtop{2}}{\@nthtop{2}}}
\DeclareRobustCommand{\thirdtop}{%
  \@ifstar{\@varnthtop{3}}{\@nthtop{3}}}
\DeclareRobustCommand{\firstbottom}{%
  \@ifstar{\@varnthbottom{1}}{\@nthbottom{1}}}
\DeclareRobustCommand{\secondbottom}{%
  \@ifstar{\@varnthbottom{2}}{\@nthbottom{2}}}
\DeclareRobustCommand{\thirdbottom}{%
  \@ifstar{\@varnthbottom{3}}{\@nthbottom{3}}}
\let\itop\firsttop
\let\iitop\secondtop
\let\iiitop\thirdtop
\let\ibottom\firstbottom
\let\iibottom\secondbottom
\let\iiibottom\thirdbottom
%    \end{macrocode}
% \end{macro}
% \iffalse
%</package>
%\fi
%
% \Finale
% \obeyspaces
% \typeout{*************************************************************}
% \typeout{*                                                           *}
% \typeout{* To finish the installation you have to move the following *}
% \typeout{* file into a directory searched by TeX:                    *}
% \typeout{*                                                           *}
% \typeout{*  actuarialsymbol.sty                                      *}
% \typeout{*                                                           *}
% \typeout{* To produce the documentation, run actuarialsymbol.dtx     *}
% \typeout{* through LaTeX.                                            *}
% \typeout{*                                                           *}
% \typeout{* Happy actuarial TeXing!                                   *}
% \typeout{*                                                           *}
% \typeout{*************************************************************}
%
% ^^A Local Variables:
% ^^A mode: doctex
% ^^A TeX-master: t
% ^^A End:
\endinput